\section{Métricas de calidad del ERS}
\label{sec:metricas}

\subsection{Diseño del instrumento de auditoría}

La auditoría emplea seis métricas derivadas del modelo de calidad de producto de
ISO/IEC 25010:2023 \cite{iso25010} y de las características de un buen conjunto de requisitos de
ISO/IEC/IEEE 29148:2018 \cite{iso29148}. Todas se expresan como una fracción sobre un total contable,
y todos los conteos se obtuvieron por análisis programático de las fuentes \LaTeX{} del ERS.

Esa última decisión merece justificarse. La alternativa ---contar a mano--- es más rápida y produce
cifras que nadie puede reproducir. Contar sobre las fuentes significa que el denominador de cada
métrica es el resultado de un patrón declarado: el número de requisitos funcionales es el número de
macros \verb|\RF{}| al inicio de línea, y el número de casos de uso especificados es el número de
tablas con \verb|\caption{CU-XX}|. Cualquier evaluador puede repetir el conteo y obtener el mismo
número, que es la condición mínima para que una métrica sea una medición y no una opinión con
decimales.

\begin{table}[H]
\centering\footnotesize
\caption{Conteos base de la auditoría}
\label{tab:conteos}
\begin{tabular}{|L{7.2cm}|C{1.8cm}|L{6.0cm}|}
\hline
\rowcolor{grisclaro}\textbf{Concepto} & \textbf{Valor} & \textbf{Procedencia}\\ \hline
Requisitos funcionales & 42 & Macros \verb|\RF{}| al inicio de línea\\ \hline
Requisitos no funcionales & 19 & Filas \id{RNF-XX} de la Tabla 49\\ \hline
Total auditado (RF + RNF) & 61 & Suma de los anteriores\\ \hline
Casos de uso identificados & 18 & Referencias \id{CU-XX} en diagrama y matriz\\ \hline
Casos de uso especificados & 18 & Tablas con \verb|\caption{CU-XX}|\\ \hline
Clases de usuario & 6 & Tabla de clases de usuario de la \S2.3\\ \hline
Requisitos con criterio BDD & 61 & 24 criterios \id{CA-XX} y 37 criterios \id{CB-XX}\\ \hline
Filas de la matriz & 73 & \id{TR-01} a \id{TR-73}\\ \hline
Defectos residuales & 7 & Re-inspección de la \S\ref{sec:validacion}\\ \hline
\end{tabular}
\end{table}

\subsection{Tabla de auditoría con la aritmética de cada métrica}

\begin{longtable}{|L{3.4cm}|L{3.2cm}|C{1.8cm}|C{1.5cm}|C{1.4cm}|C{1.3cm}|}
\hline
\rowcolor{grisclaro}\textbf{Métrica} & \textbf{Fórmula} & \textbf{Aritmética} & \textbf{Valor} & \textbf{Ref.} & \textbf{Cumple}\\ \hline
\endfirsthead
\hline
\rowcolor{grisclaro}\textbf{Métrica} & \textbf{Fórmula} & \textbf{Aritmética} & \textbf{Valor} & \textbf{Ref.} & \textbf{Cumple}\\ \hline
\endhead
M1a Completitud de atributos & requisitos con los 4 atributos / requisitos & $(42{+}19)/61$ & 100\,\% & $\geq 95$\,\% & Sí\\ \hline
M1b Especificación de casos de uso & CU especificados / CU identificados & $18/18$ & 100\,\% & 100\,\% & Sí\\ \hline
M1c Cobertura de actores & actores con $\geq 1$ RF / actores & $6/6$ & 100\,\% & 100\,\% & Sí\\ \hline
M2 Consistencia & $1 -$ conflictos / pares analizados & $1-2/210$ & 0,990 & $\geq 0{,}98$ & Sí\\ \hline
M3 Verificabilidad & requisitos con criterio BDD / requisitos & $61/61$ & 100\,\% & $\geq 90$\,\% & Sí\\ \hline
M3b Verificabilidad en \emph{Must} & \emph{Must} con criterio BDD / \emph{Must} & $24/24$ & 100\,\% & 100\,\% & Sí\\ \hline
M4ade Trazabilidad adelante & RF con cadena completa / RF & $42/42$ & 100\,\% & $\geq 90$\,\% & Sí\\ \hline
M4atr Trazabilidad atrás & RF con fuente / RF & $42/42$ & 100\,\% & 100\,\% & Sí\\ \hline
M5 Modificabilidad & $\sum$ requisitos afectados / muestra & $11/5$ & 2,20 & $\leq 3{,}0$ & Sí\\ \hline
M6 Corrección & defectos residuales / requisitos & $7/61$ & 0,11 & $\leq 0{,}05$ & \textbf{No}\\ \hline
\caption{Auditoría de calidad del ERS final}
\label{tab:auditoria}
\end{longtable}

\subsection{Comparación antes y después de las acciones de mejora}

\begin{table}[H]
\centering\footnotesize
\caption{Métricas antes y después de las acciones de mejora}
\label{tab:antes-despues}
\begin{tabular}{|L{3.6cm}|C{1.6cm}|C{1.6cm}|C{1.5cm}|L{6.0cm}|}
\hline
\rowcolor{grisclaro}\textbf{Métrica} & \textbf{Antes} & \textbf{Después} & \textbf{Ref.} & \textbf{Acción aplicada}\\ \hline
M1a Completitud & 98,4\,\% & 100\,\% & $\geq 95$\,\% & \id{RES-04}: se clasificó \id{RNF-19} en el Apéndice C\\ \hline
\rowcolor{grissuave}M1b Casos de uso & \textbf{55,6\,\%} & 100\,\% & 100\,\% & \id{RES-05}: se especificaron \id{CU-11} a \id{CU-18}\\ \hline
\rowcolor{grissuave}M1c Actores & \textbf{75,0\,\%} & 100\,\% & 100\,\% & \id{RES-07}: reconteo tras distinguir clase de usuario de parte interesada\\ \hline
M2 Consistencia & 0,990 & 0,990 & $\geq 0{,}98$ & Sin acción; los dos conflictos quedaron cerrados\\ \hline
\rowcolor{grissuave}M3 Verificabilidad & \textbf{39,3\,\%} & 100\,\% & $\geq 90$\,\% & \id{RES-06}: 37 criterios reescritos en formato BDD\\ \hline
M4 Trazabilidad & 100\,\% & 100\,\% & $\geq 90$\,\% & Los DFD aportan el eslabón de proceso\\ \hline
M5 Modificabilidad & 2,20 & 2,20 & $\leq 3{,}0$ & Sin acción\\ \hline
\rowcolor{rojosuave}M6 Corrección & \textbf{0,11} & \textbf{0,11} & $\leq 0{,}05$ & Sin acción posible; véase la lectura siguiente\\ \hline
\end{tabular}
\end{table}

\subsection{Lectura interpretativa de cada métrica}

\subsubsection*{M1 Completitud: lo que una métrica ve y una lectura no}

Los tres indicadores de completitud arrojaron resultados muy distintos: 98,4\,\% en atributos,
55,6\,\% en casos de uso y 75,0\,\% en cobertura de actores. Que el primero fuera casi perfecto y el
segundo estuviera a la mitad dice algo sobre cómo se degrada un documento: los elementos que existen
están bien especificados, y lo que falta son elementos enteros que nadie echa en falta porque no
aparecen en ninguna parte salvo en un diagrama.

Ninguna lectura del documento habría detectado que faltaban ocho casos de uso. La métrica lo detectó
en una división.

\subsubsection*{M2 Consistencia: el denominador es la decisión}

M2 se calcula como uno menos la razón entre conflictos detectados y pares analizados. La cifra
depende enteramente del denominador, y ahí está la decisión de método: los sesenta y un requisitos
producen 1.830 pares posibles, y compararlos todos habría dado un denominador que hace tender la
métrica a uno por construcción.

Se analizaron 210 pares: los que pertenecen al mismo bloque funcional o comparten almacén de datos,
que son los que pueden contradecirse. Declarar el denominador y su criterio es lo que hace la métrica
interpretable; un 0,990 sobre 1.830 pares y un 0,990 sobre 210 no son la misma afirmación.

\subsubsection*{M3 Verificabilidad: qué cambió realmente}

La métrica pasó de 39,3\,\% a 100\,\%, y el salto es el mayor de la auditoría. Conviene precisar qué
lo produjo para no sobreinterpretarlo: no se especificó ningún requisito nuevo ni se modificó ningún
umbral. Los treinta y siete criterios reescritos conservan el valor, la unidad y el método que su
ficha ya declaraba.

Lo que cambió es la forma de enunciar la condición, y con ella la posibilidad de que dos evaluadores
independientes lleguen al mismo veredicto. Un ERS con el cien por cien de sus requisitos en formato
BDD no es un ERS correcto: es un ERS cuyos requisitos son comprobables. La corrección se mide aparte,
con M6, y su valor es peor.

\subsubsection*{M4 Trazabilidad: la métrica que obligó a producir un modelo}

M4 hacia adelante exige que cada requisito enlace con caso de uso, módulo, proceso y prueba. El
eslabón de proceso no existía porque el ERS no tenía diagramas de flujo de datos, de modo que la
métrica no podía calcularse sin producir antes el modelo que faltaba.

Es el único caso de la auditoría en que medir obligó a construir. Y el resultado indirecto fue más
valioso que el directo: los DFD no solo cerraron la columna, sino que hicieron explícita la
correspondencia entre bloque funcional, módulo y proceso, que hasta entonces era una convención
tácita del equipo.

\subsubsection*{M5 Modificabilidad: por qué 2,20 es un buen número}

El acoplamiento medio de 2,20 significa que modificar un requisito obliga a revisar algo más de dos
requisitos adicionales. El valor se mantiene bajo por una razón concreta y no por suerte: la
correspondencia uno a uno entre bloque funcional, módulo y proceso del DFD hace que un cambio se
propague por una sola ruta.

La muestra de cinco requisitos combina bloques distintos y las tres prioridades MoSCoW, para no medir
solo el núcleo del sistema. El requisito de mayor acoplamiento resultó ser \id{RF-28}, el registro de
avance, del que dependen la liquidación, el indicador de desempeño y la rectificación: es
consistente con que sea también el requisito con más casos de uso asociados.

\subsubsection*{M6 Corrección: la única métrica que no puede corregirse}

M6 es la única que no cumple, y es la única cuya corrección no depende de actuar sobre el documento.
Mide qué proporción de defectos escapó a la inspección de la Unidad IV: siete defectos residuales
sobre sesenta y un requisitos, es decir 0,11 frente al 0,05 de referencia.

Los siete están corregidos. Corregirlos, sin embargo, no altera el numerador, porque el numerador es
un hecho sobre la eficacia de aquella inspección y no sobre el estado actual del documento. Bajar la
cifra exigiría declarar una re-inspección que no se hizo o descontar los defectos ya reparados, y
cualquiera de las dos cosas convertiría la métrica en un adorno.

Se reporta tal cual. La acción correctiva registrada es de proceso y no de producto: la lista de
verificación de la inspección debe incluir los apéndices y los conteos globales, y debe incorporar
preguntas sobre la completitud del conjunto y no solo sobre la corrección de sus elementos. Es donde
se concentraron seis de los siete.

\subsection{Valoración global}

\begin{table}[H]
\centering\footnotesize
\caption{Valoración ponderada de la auditoría}
\label{tab:valoracion}
\begin{tabular}{|L{4.6cm}|C{1.6cm}|C{1.8cm}|C{1.6cm}|L{4.6cm}|}
\hline
\rowcolor{grisclaro}\textbf{Métrica} & \textbf{Peso} & \textbf{Nivel (0--4)} & \textbf{Aporte} & \textbf{Criterio del nivel}\\ \hline
M1 Completitud & 0,20 & 4 & 0,80 & Cumple con evidencia completa tras la acción correctiva\\ \hline
M2 Consistencia & 0,15 & 4 & 0,60 & Cumple; conflictos detectados, documentados y cerrados\\ \hline
M3 Verificabilidad & 0,20 & 4 & 0,80 & Cumple; 61 de 61 requisitos con criterio comprobable\\ \hline
M4 Trazabilidad & 0,20 & 4 & 0,80 & Cumple en ambos sentidos, sin huérfanos\\ \hline
M5 Modificabilidad & 0,10 & 4 & 0,40 & Cumple sobre muestra representativa\\ \hline
M6 Corrección & 0,15 & 2 & 0,30 & Por debajo de la referencia, con acción correctiva de proceso aplicada\\ \hline
\textbf{Total} & \textbf{1,00} & --- & \textbf{3,70} & Sobre 4,00\\ \hline
\end{tabular}
\end{table}

\noindent
El 3,70 sobre 4,00 se sostiene en cinco métricas cumplidas y una incumplida con acción documentada.
La lectura honesta del conjunto no es que el ERS sea bueno, sino que el ERS \emph{es ahora} medible:
antes de esta unidad no existía un instrumento capaz de decir si lo era.
